\documentclass[a4paper]{article}

\usepackage[fontset=fandol]{ctex}
\usepackage{amsmath, amssymb}
\usepackage{graphicx}
\usepackage[margin=2.5cm]{geometry}
\usepackage[most]{tcolorbox}
\usepackage{etoolbox}
\usepackage{listings}
\usepackage{hyperref}
\usepackage{booktabs}
\usepackage{subcaption}
\usepackage{float}
\usepackage{tikz}
\IfFileExists{pgfplots.sty}{
    \usepackage{pgfplots}
    \pgfplotsset{compat=1.18}
}{}
\hypersetup{hidelinks}

\newtcolorbox{knowledgebox}[1]{
    enhanced, colback=blue!5!white, colframe=blue!75!black, colbacktitle=blue!75!black,
    coltitle=white, fonttitle=\bfseries, title=#1, attach boxed title to top left={yshift=-2mm, xshift=2mm},
    boxrule=1pt, sharp corners
}

\newtcolorbox{importantbox}[1]{
    enhanced, colback=yellow!10!white, colframe=yellow!80!black, colbacktitle=yellow!80!black,
    coltitle=black, fonttitle=\bfseries, title=#1, sharp corners
}

\newtcolorbox{warningbox}[1]{
    enhanced, colback=red!5!white, colframe=red!75!black, colbacktitle=red!75!black,
    coltitle=white, fonttitle=\bfseries, title=#1, sharp corners
}

\lstset{
    language=Python,
    basicstyle=\ttfamily\small,
    keywordstyle=\color{blue},
    stringstyle=\color{red!60!black},
    commentstyle=\color{green!60!black},
    breaklines=true,
    frame=single,
    numbers=left,
    numberstyle=\tiny\color{gray},
    captionpos=b,
    extendedchars=false
}

\newcommand{\notetitle}{CS224R 第 14 讲：Exploration \& Meta-Exploration}
\newcommand{\noteauthors}{Codex 基于 Stanford CS224R 2025 春季课程视频、字幕与官方 slides 重写}
\newcommand{\notedate}{2026-04-13}
\newcommand{\videochannel}{Stanford Online}
\newcommand{\videopublishdate}{2025年12月08日}
\newcommand{\videoduration}{1:12:42}
\newcommand{\videourl}{https://www.youtube.com/watch?v=4tlSKdi8teU}
\newcommand{\videocoverpath}{cover.jpg}

\begin{document}

\begin{titlepage}
\centering
{\Large 课程笔记\par}
\vspace{1.2cm}
{\huge\bfseries \notetitle\par}
\vspace{0.8cm}
{\large \noteauthors\par}
\vspace{0.3cm}
{\large \notedate\par}
\vspace{1.2cm}

\includegraphics[width=0.82\textwidth,height=0.45\textheight,keepaspectratio]{\videocoverpath}\par

\vfill
\begin{tcolorbox}[width=0.9\textwidth, colback=black!2!white, colframe=black!60, sharp corners]
\textbf{视频作者/频道}：\videochannel\par
\textbf{发布日期}：\videopublishdate\par
\textbf{视频时长}：\videoduration\par
\textbf{视频链接}：\href{\videourl}{\nolinkurl{\videourl}}
\end{tcolorbox}
\end{titlepage}

\tableofcontents
\newpage

\section{为什么探索如此困难}
这一讲开篇没有直接讲算法，而是先建立“探索问题到底难在哪里”的直觉。讲者用 Atari 里的 \textit{Breakout} 与 \textit{Montezuma's Revenge} 对比说明：在前者中，奖励频繁且行为-结果关系直观，标准 RL 往往能学得不错；而在后者中，关键奖励事件稀疏、延迟长、局部动作并不直接显出价值，因此“从零发现高回报行为链条”非常困难。

\begin{figure}[H]
\centering
\includegraphics[width=0.9\textwidth]{pdf_pages/page-05.png}
\caption{Montezuma's Revenge 例子说明：拿钥匙、开门等关键事件都很稀疏，算法并不知道哪些 sprite 值得关注。}
\end{figure}

课程随后引入 Mao 这个“规则只能靠试错发现”的牌类游戏，让学生站到算法视角去想：如果你只知道“犯错会受罚”，却不知道规则是什么，也不知道惩罚的语义是否合理，那么任何长时程任务都会变得异常艰难。任务越长、规则越隐蔽、奖励越稀疏，探索就越难。

这使得“探索”与“利用”看起来像两个定义：
\begin{itemize}
\item 探索：尝试尚未验证的新行为，希望发现更高回报；
\item 利用：执行已知回报最高的行为。
\end{itemize}
但课程强调，这其实是同一个问题的两面。真正的挑战不是简单地“随机试更多动作”，而是在有限交互预算下，决定何时继续沿用当前最好策略、何时为未来收益而承担短期损失。

\begin{importantbox}{本讲的基本立场}
探索问题的本质，不是“怎样更随机”，而是“怎样在不确定性下为未来信息价值付费”。
\end{importantbox}

为了让直觉更接近日常，讲者给出三个典型例子：选餐馆、广告投放、油井勘探。这些例子都说明，一旦长期收益取决于今天是否愿意尝试新选项，探索与利用就不可分割。

\subsection{本章小结}
探索之所以难，不是因为“找不到随机性”，而是因为高回报往往需要长链条、稀疏反馈和对环境规则的推断。稀疏奖励只是表象，更深层的问题是不确定性与长时程信用分配。

\section{Bandit 视角：如何定义“好的探索”}
为了让讨论可分析，课程先退到最简单的 setting：multi-armed bandit。它只有单步 horizon、没有状态，但奖励是随机的，因此非常适合分析探索策略的优劣。

\begin{figure}[H]
\centering
\includegraphics[width=0.86\textwidth]{pdf_pages/page-09.png}
\caption{从 bandit 到大规模无结构 MDP，探索问题会迅速从可分析走向理论与计算上的困难。}
\end{figure}

在 bandit 里，课程引入了后续最重要的评价指标：regret。直观上，它衡量的是“如果我一开始就知道最佳动作，和我实际所选动作相比，累计少拿了多少回报”。

\[
\mathrm{Regret}(T) = T \max_{a} \mathbb{E}[r(a)] - \sum_{t=1}^{T} r(a_t).
\]
这里：
\begin{itemize}
\item $T$ 是累计交互轮数；
\item $a$ 表示任意候选动作；
\item $\mathbb{E}[r(a)]$ 是动作 $a$ 的期望回报；
\item $a_t$ 是第 $t$ 轮实际采取的动作；
\item $\mathrm{Regret}(T)$ 越小，表示探索-利用权衡越好。
\end{itemize}

这个定义的关键之处在于，它不是只看最终找到没找到最优臂，而是把整个学习过程都纳入评价。一个算法如果一直随机选动作，regret 会近似线性增长；一个有效算法则希望 regret 以更慢的速度增长。

\begin{figure}[H]
\centering

\includegraphics[width=0.92\textwidth]{pdf_pages/page-11.png}
\caption{Bandit 中两条经典思路：UCB 的“在不确定时乐观”与 posterior sampling 的“按后验概率匹配”。}
\end{figure}

课程随后并列介绍两类经典策略：
\begin{itemize}
\item UCB：对均值估计再加上不确定性 bonus，优先尝试“可能比现在看起来更好”的臂。
\item posterior sampling：维护每个臂的后验分布，从中采样一个世界模型，再在该模型下选当前最优动作。
\end{itemize}

UCB 可以写成
\[
a_t=\arg\max_{a}\left(\bar{r}_a + c\,\sigma_a\right).
\]
这里：
\begin{itemize}
\item $\bar{r}_a$ 是动作 $a$ 的当前平均回报估计；
\item $\sigma_a$ 是该动作的不确定性或方差估计；
\item $c$ 控制探索 bonus 的强度；
\item $a_t$ 是当前轮被选中的动作。
\end{itemize}

这个公式浓缩了 UCB 的全部直觉：未知并不直接等于好，但未知值得被加分，因为它可能隐藏更高收益。

\begin{knowledgebox}{为什么 bandit 重要}
它不是因为现实问题都像 bandit 那样简单，而是因为这里能清楚定义 regret、比较算法、分析长期收益，从而为更复杂的探索问题建立语言和直觉。
\end{knowledgebox}

\subsection{本章小结}
Bandit 提供了探索问题最干净的实验台。通过 regret，我们不再只关心“最终找没找到最优解”，而是开始关心整个学习过程的累计损失；通过 UCB 与 posterior sampling，我们得到两种最经典的探索思想。

\section{从 Bandit 到机器人与大模型}
讲者紧接着提醒：\textbf{不要误以为 bandit 里的优雅结论会自动延伸到大规模 MDP。} 一旦有状态、有长时程依赖、奖励稀疏且环境没有已知结构，探索很快就变成理论上和计算上都极难的问题。

课程在这一部分给了一个非常务实的结论。尽管学界尝试把 UCB、posterior sampling 等思想推广到机器人控制等大 MDP 中，但“从零探索”在现实大型问题里通常是不现实的。因此工业和前沿研究更常采用三类补救：
\begin{itemize}
\item 借助 demonstrations 或在 demonstrations 上预训练的基础模型；
\item 尽量使用 shaped rewards；
\item 让模型依赖更强的先验知识，而不是从纯随机交互中发现一切。
\end{itemize}

这也是为什么讲者会把机器人和 LLM 并列提起。大型语言模型之所以看起来“探索能力很强”，并不是因为它们在线上从零开始探索得更好，而是因为海量预训练已经把世界知识、语义先验和动作模式压进了参数里。某种意义上，它们是在用预训练替代在线探索。

\subsection{本章小结}
Bandit 给我们方法论，但现实中的机器人和 LLM 更多依赖先验、数据与奖励工程。真正困难的大规模探索问题，很少靠“从零试出来”解决。

\section{测试时探索与 Meta-Exploration}
在前半部分建立了探索基本图景之后，课程把话题拉回 Meta-RL。这里的问题变成：如果测试时代理必须先探索，再利用探索结果执行任务，那么该如何训练这种行为？

讲者先指出一个非常常见的 baseline：端到端优化探索 episode 与执行 episode 的总回报。这种做法原则上最直接，因为它直接对最终目标负责。

\begin{figure}[H]
\centering
\includegraphics[width=0.9\textwidth]{pdf_pages/page-16.png}
\caption{元探索中的耦合问题：探索决定能否收集信息，执行决定信息能否转化为回报，两者相互依赖。}
\end{figure}

但它也有一个经典 failure mode：coupling problem。若代理不会探索，它就收集不到执行所需的信息；若它不会执行，即使探索对了也拿不到高回报。课程用“在新厨房先找食材再做饭”的例子说明，这会导致极差的样本效率和糟糕的局部最优。

这时可以回到上一讲提到的另一条路线：posterior sampling。PEARL 就是在元学习框架里显式建模任务潜变量后验，再按当前采样到的任务假设去行动。这种做法虽然不再是纯粹黑盒端到端，但它把“任务推断”和“策略执行”分成了两个更清晰的模块。

课程随后介绍 DREAM。它的核心思想不是直接逼网络同时学好探索与执行，而是先学会抽取\textbf{与任务有关的信息表示}，并通过 bottleneck 压掉无关信息，从而把探索与执行部分解耦。

\begin{figure}[H]
\centering
\includegraphics[width=0.9\textwidth]{pdf_pages/page-20.png}
\caption{DREAM 的关键想法：先学习能识别任务差异的瓶颈表示，再让执行策略围绕该表示工作。}
\end{figure}

从 lecture 的叙述看，DREAM 的价值在于两点：
\begin{itemize}
\item 它保留了“测试时需要探索”的场景，而不是偷换成显式任务标签。
\item 它把信息收集问题转化为“哪些信息对任务真正 relevant”的表示学习问题。
\end{itemize}

讲者还强调，DREAM 在理论上能在较弱假设下与端到端目标保持一致，在实践上则通常比完全耦合的端到端方法更高效。

\begin{figure}[H]
\centering
\includegraphics[width=0.88\textwidth]{pdf_pages/page-25.png}
\caption{在稀疏奖励 3D 导航问题上，DREAM 类方法需要先绕路读取线索，再去目标物体位置。}
\end{figure}

\begin{warningbox}{本讲的研究判断}
元探索真正难的不是“有没有任务记忆”，而是“如何获得足以区分任务、又不和执行学习彼此拖累的信息”。这也是为什么表示学习和解耦成为关键。
\end{warningbox}

\subsection{本章小结}
测试时探索把普通探索问题推进了一步：不仅要发现高回报动作，还要先发现“这个任务到底是什么”。端到端训练虽最完整，却很容易耦合失效；PEARL 与 DREAM 则尝试通过结构化后验或信息瓶颈来缓解这一点。

\section{Meta-RL 在教育中的应用}
本讲最后给了一个非常具体也很有启发的应用：把学生程序看作不同任务，用 Meta-RL 代理在程序里“探索 bug”，再据此生成反馈和辅助评分。

\begin{figure}[H]
\centering
\includegraphics[width=0.88\textwidth]{pdf_pages/page-30.png}
\caption{课程把学生程序视作不同任务：代理先交互探索程序行为，再定位 bug 并辅助反馈或评分。}
\end{figure}

Code.org 的 Bounce 和 Stanford CS106A 的 Breakout 都是很好的例子。不同学生提交的程序会有不同 bug，有的球会被球拍“串住”，有的会无限生成物体。教师或助教的工作，本质上是通过交互去发现“这个程序的潜在规则和故障模式是什么”。这与 Meta-RL 的任务适应范式非常接近。

课程报告了两个很直接的收益：
\begin{itemize}
\item 在 Bounce 任务中，系统能学会有针对性的探索行为，主动触发关键事件来暴露 bug。
\item 在 CS106A 2023 春季学期的评分试验中，辅助系统带来了约 $44\%$ 更快、约 $6\%$ 更准确的 grading。
\end{itemize}

这部分的价值不只在于结果数字，而在于它说明 Meta-RL 并非只能用于导航、机器人或 bandit。只要存在“先探索行为，再据此推断任务结构”的过程，Meta-RL 就可能成为一个合适的建模工具。

\subsection{本章小结}
教育应用把本讲的抽象问题落到了非常具体的交互系统上。它提醒我们：探索不仅是找奖励，也可以是为了更快地发现 bug、生成反馈和支持人类决策。

\section{总结与延伸}
这一讲分成两条主线。第一条主线是探索本身：为什么稀疏奖励、长时程任务和规则不透明会让探索变难；为什么 bandit 是研究探索的好起点；为什么 UCB 与 posterior sampling 代表了两种经典范式。第二条主线是 meta-exploration：如果测试时必须先探索当前任务，再利用收集到的信息，那么训练会遇到严重的耦合问题。

在研究判断上，这讲非常克制。它没有声称“找到一种万能探索算法”，反而一直在强调 problem structure 的重要性：
\begin{itemize}
\item 简单结构问题可以谈 regret 与近最优算法；
\item 复杂大 MDP 往往依赖 demonstrations、预训练和 shaped rewards；
\item 元探索则需要额外考虑任务识别、信息瓶颈和探索-执行解耦。
\end{itemize}

从课程脉络看，这讲实际上起到了桥梁作用。它一边把上一讲的 Meta-RL 进一步推到“测试时必须探索”的 setting，一边又为下一讲关于 autonomy、long-horizon task 和 hierarchy 埋下伏笔。对读者而言，最值得带走的不是某一个算法名字，而是这句隐含的工程判断：\textbf{探索难点的本质，常常取决于你缺的到底是动作尝试、任务识别，还是足够强的先验。}

\end{document}
